\documentclass[12pt, a4paper]{article}

\author{Lucas Sernik}
\date{\today}

\begin{document}
{\large Lucas Sernik's Peer Review of Raja Chinnakotla First Draft}

\vspace{1cm}

\quad Raja Chinnakotla wants to explore how the linguistical difference in Gujarat, a state in India, contributed to the intercommunal violence in the year 2002. Even though there are different theories to explain why the violence occured continuously during the period of three days, none of them present a theory fundamented on empirical data and modelling with some media component.

\quad Raja will try to use data from the Indian 2001 census, which correlates population in a given town/district with their religion, languages spoke, distance to nearest train stations, and other related control variables. He will then use a model to evaluate if there is any relationship between neighborhoods/districts with a Gujarat Samachar headquarter with the number of people killed and/or the number of buildings destroyed.

\quad I think that the overall theme chosen is very interesting. I was not aware of the Gujarat 2002 riots and its ethnical impacts in India, and I an glad that, if the government is/will try to erase part of the history, there are peope who are concerned with documenting what happened before it is too late.

\quad Discussing more about the economic aspect of the paper, I suggest that you take a look into other countries in Africa and Latin America, looking for similar conflicts between two different ethnic groups. These might help you with some citations and give more empirical evidence that might support your claims. Also, you might get some insight on the methodologies used and any endogeneity discussions, which will help you frame your own approach better

\quad More specifically, Brazil, Bolivia, Colombia, and other countries in Central America should have some interesting papers exposing events similar to the riot, as they have a substantial indigenous population, and miscegenation, with several cultures developing alongside each other.

\quad Moreover, some papers analyzing these regions might have relevant economic data supporting the fact that violence hinders economic growth. Brazil contains several favela's (slums) controlled by organized crime, with some neighborhoods being disputed by different factions, where we could apply the theory that violent regions are more costly for firms to set up and operate.

\quad Related to the media, you might also find interesting to compare media outlets in the same language but with different political views.

\quad Another way to explore the media inciting violence are the recent protests in Nepal. Since they were led by Gen Z through media, you could make an argument that younger people had their own media format and their versions of what was happening (analogous to their language) while older generations were still informed by more traditional media outlets, which led to the conflict (I don't know if any part of the Nepal argument is actually true, just brainstorming some events that might help you).
\end{document}